\documentclass[11pt,a4paper]{article}

% ---------- Codificación y tipografía ----------
\usepackage[utf8]{inputenc}
\usepackage[T1]{fontenc}
\usepackage{lmodern}
\usepackage[spanish,es-nodecimaldot]{babel}
\usepackage{microtype}
\usepackage{amsmath,amssymb}
\usepackage{geometry}
\geometry{top=2.3cm,bottom=2.3cm,left=2.2cm,right=2.2cm}
\usepackage{enumitem}
\usepackage{booktabs}
\usepackage{longtable}
\usepackage{array}
\usepackage{fancyhdr}
\usepackage{titlesec}
\usepackage{setspace}
\usepackage{multicol}
\usepackage{needspace}

\setstretch{1.08}
\setlength{\parindent}{0pt}
\setlength{\parskip}{4pt}

% ---------- Encabezado y pie ----------
\pagestyle{fancy}
\fancyhf{}
\lhead{\small Banco de Preguntas PEU --- SUNASS}
\rhead{\small Módulo II: Teoría del Productor}
\cfoot{\small \thepage}
\renewcommand{\headrulewidth}{0.4pt}
\renewcommand{\footrulewidth}{0pt}

% ---------- Títulos ----------
\titleformat{\section}{\large\bfseries}{\thesection}{0.8em}{}
\titlespacing*{\section}{0pt}{14pt}{6pt}
\titleformat{\subsection}{\normalsize\bfseries}{\thesubsection}{0.6em}{}
\titlespacing*{\subsection}{0pt}{10pt}{4pt}

% ---------- Comandos del examen ----------
\newcounter{qnum}
\newcommand{\q}[1]{%
  \needspace{8\baselineskip}%
  \par\medskip\refstepcounter{qnum}%
  \noindent\textbf{\theqnum.}\ #1\par\nopagebreak[4]
}
\newenvironment{opts}
  {\begin{enumerate}[label=\Alph*),leftmargin=2.0em,itemsep=1pt,topsep=2pt,parsep=0pt,partopsep=0pt]}
  {\end{enumerate}}
\makeatletter
\@itempenalty=10000
\makeatother
\newcommand{\sol}[3]{\par\noindent\textbf{#1. #2.}\ #3\par\smallskip}

\begin{document}

% ============================================================
% PORTADA
% ============================================================
\begin{titlepage}
\thispagestyle{empty}
\centering
\vspace*{1.2cm}

{\Large\bfseries BANCO DE PREGUNTAS}\\[0.35cm]
{\large EXAMEN DE ADMISIÓN}\\[0.15cm]
{\large PROGRAMA DE EXTENSIÓN UNIVERSITARIA (PEU)}\\[0.15cm]
{\large SUPERINTENDENCIA NACIONAL DE SERVICIOS DE SANEAMIENTO}\\[0.05cm]
{\large SUNASS}\\[1.0cm]

\rule{\linewidth}{0.9pt}\\[0.55cm]
{\LARGE\bfseries MÓDULO II: TEORÍA DEL PRODUCTOR}\\[0.35cm]
{\large Especialidad: Economía}\\[0.55cm]
\rule{\linewidth}{0.9pt}\\[1.0cm]

\begin{tabular}{ll}
\toprule
Número de preguntas: & 90 \\
Bloque A: & 30 preguntas \\
Bloque B: & 35 preguntas \\
Bloque C: & 25 preguntas \\
Duración sugerida: & 150 minutos \\
Nivel: & Microeconomía intermedia--avanzada \\
Formato: & Opción múltiple (A, B, C, D) \\
\bottomrule
\end{tabular}

\vfill
{\small Material de preparación académica. Uso interno.}\\[0.15cm]
{\small Edición 2026}
\end{titlepage}

% ============================================================
% INSTRUCCIONES
% ============================================================
\section*{Instrucciones generales}

\begin{enumerate}[label=\arabic*.,leftmargin=1.6em,itemsep=3pt]
  \item El presente banco contiene \textbf{90 preguntas de opción múltiple}. Cada pregunta tiene una sola respuesta correcta.
  \item No se descuenta por respuesta incorrecta; se recomienda responder todas las preguntas.
  \item El tiempo sugerido de resolución es de \textbf{150 minutos}. Use calculadora no programable.
  \item Las preguntas están organizadas en tres bloques: fundamentos teóricos, aplicaciones cuantitativas y casos regulatorios del sector saneamiento.
  \item El \textbf{solucionario completo} y la \textbf{tabla de claves rápidas} se encuentran al final. Se recomienda resolver el examen antes de consultarlos.
\end{enumerate}

\section*{Estructura del banco}

\begin{center}
\renewcommand{\arraystretch}{1.25}
\begin{tabular}{@{}p{1.4cm}>{\raggedright\arraybackslash}p{7.6cm}cc@{}}
\toprule
\textbf{Bloque} & \textbf{Contenido} & \textbf{Preguntas} & \textbf{Peso} \\
\midrule
A & Fundamentos teóricos y conceptuales: tecnología, isocuantas, rendimientos, isocostos, minimización de costos, dualidad, competencia y bienestar del productor. & 1--30 & 33.3\% \\
B & Aplicaciones cuantitativas y optimización: Cobb-Douglas, CES, Leontief, sustitutos perfectos, productividades, funciones de costo, equilibrio competitivo y economías de escala. & 31--65 & 38.9\% \\
C & Casos con implicancias tarifarias y regulatorias del sector saneamiento: costos medios y marginales, monopolio natural, RPI-X, tasa de retorno, ANF, productividad y asignación de costos. & 66--90 & 27.8\% \\
\bottomrule
\end{tabular}
\end{center}

\section*{Competencias evaluadas}

\noindent El banco evalúa la capacidad de interpretar la tecnología de producción, resolver problemas de minimización de costos y maximización de beneficios, derivar funciones de costo y oferta, distinguir economías de escala y de alcance, y aplicar estos conceptos al diseño y evaluación de tarifas de agua potable y alcantarillado.

\section*{Notación principal}

{\small
\renewcommand{\arraystretch}{1.0}
\begin{center}
\begin{tabular}{@{}ll@{\hspace{1.0cm}}ll@{}}
\toprule
$Q$ & Producción & $C(Q)$ & Costo total \\
$K,L$ & Capital y trabajo & $CMg$ & Costo marginal \\
$w,r$ & Precios del trabajo y del capital & $CMe$ & Costo medio \\
$PMg$ & Producto marginal & $CV$ & Costo variable \\
$PMe$ & Producto medio & $CF$ & Costo fijo \\
$RMST$ & Relación marginal de sustitución técnica & $IT$ & Ingreso total \\
$\sigma$ & Elasticidad de sustitución & $\pi$ & Beneficio económico \\
\bottomrule
\end{tabular}
\end{center}
}

\newpage
\setcounter{qnum}{0}

\section*{Bloque A. Fundamentos teóricos y conceptuales}

\noindent\textit{Preguntas 1 a 30. Evalúan la comprensión analítica de la tecnología de producción, las isocuantas, los rendimientos a escala, la minimización de costos, la dualidad, la oferta competitiva y los conceptos de beneficio y excedente del productor.}

\q{La función de producción $Q=f(K,L)$ representa:}
\begin{opts}
    \item el costo mínimo necesario para producir una cantidad de producto.
    \item el máximo producto factible con insumos dados y tecnología vigente.
    \item el beneficio económico máximo alcanzable en equilibrio de mercado.
    \item la cantidad mínima de insumos requerida para evitar pérdidas.
\end{opts}

\q{Una isocuanta es el conjunto de combinaciones de insumos que:}
\begin{opts}
    \item generan exactamente el mismo costo total de producción.
    \item reportan exactamente el mismo ingreso total de ventas.
    \item permiten alcanzar exactamente el mismo nivel de producto.
    \item producen exactamente el mismo beneficio económico neto.
\end{opts}

\q{La relación marginal de sustitución técnica (RMST) entre trabajo y capital se define como:}
\begin{opts}
    \item capital sustituible por una unidad más de trabajo manteniendo $Q$.
    \item variación total del producto ante cambios simultáneos en insumos.
    \item pendiente constante de la recta presupuestaria o recta isocosto.
    \item cociente directo entre el salario nominal y la renta del capital.
\end{opts}

\q{Si $PMg_L$ y $PMg_K$ son los productos marginales del trabajo y del capital, la RMST de trabajo por capital es:}
\begin{opts}
    \item $w/r$.
    \item $r/w$.
    \item $PMg_L/PMg_K$.
    \item $PMg_K/PMg_L$.
\end{opts}

\q{Una RMST decreciente a lo largo de una isocuanta implica que las isocuantas son:}
\begin{opts}
    \item líneas rectas.
    \item cóncavas respecto del origen.
    \item convexas respecto del origen.
    \item en forma de L.
\end{opts}

\q{Existen rendimientos constantes a escala cuando, al duplicar todos los insumos:}
\begin{opts}
    \item el producto aumenta menos del doble.
    \item el producto aumenta más del doble.
    \item el producto permanece constante.
    \item el producto se duplica exactamente.
\end{opts}

\q{Para la función Cobb-Douglas $Q=AK^\alpha L^\beta$, los rendimientos a escala están dados por:}
\begin{opts}
    \item $\alpha/\beta$.
    \item $\alpha+\beta$.
    \item $\alpha\beta$.
    \item $\max(\alpha,\beta)$.
\end{opts}

\q{Si en una función Cobb-Douglas $Q=AK^\alpha L^\beta$ se cumple $\alpha+\beta>1$, la tecnología presenta rendimientos:}
\begin{opts}
    \item decrecientes a escala.
    \item crecientes a escala.
    \item constantes a escala.
    \item negativos a escala.
\end{opts}

\q{La ley de los rendimientos marginales decrecientes establece que:}
\begin{opts}
    \item en el largo plazo todos los insumos productivos son variables.
    \item al duplicar ambos insumos el volumen de producción se duplica.
    \item la producción conjunta resulta siempre más barata que separada.
    \item con un factor fijo, el $PMg$ del factor variable acaba cayendo.
\end{opts}

\q{La relación entre producto medio ($PMe$) y producto marginal ($PMg$) implica que:}
\begin{opts}
    \item el $PMe$ alcanza su máximo cuando el $PMg$ es máximo.
    \item $PMg=PMe$ cuando el $PMe$ es cero.
    \item $PMg=PMe$ en el máximo del $PMe$.
    \item el $PMg$ es siempre mayor que el $PMe$.
\end{opts}

\q{La etapa II de la función de producción, considerada la región racional de producción, comprende:}
\begin{opts}
    \item desde el origen hasta el máximo del producto medio.
    \item desde el máximo del producto medio hasta que el producto marginal se hace cero.
    \item desde que el producto marginal se hace cero hasta que se vuelve negativo.
    \item toda la función de producción.
\end{opts}

\q{La recta isocosto $C=wL+rK$ tiene una pendiente igual a:}
\begin{opts}
    \item $-r/w$.
    \item $-w/r$.
    \item $-PMg_L/PMg_K$.
    \item $-P/Q$.
\end{opts}

\q{En el óptimo de minimización de costos para una producción dada, se cumple:}
\begin{opts}
    \item $RMST_{LK}=w/r$.
    \item $RMST_{LK}=r/w$.
    \item $PMg_L=PMg_K$.
    \item $wL=rK$.
\end{opts}

\q{La senda de expansión de la empresa es:}
\begin{opts}
    \item la ruta de combinaciones óptimas de insumos al variar $Q$ a precios fijos.
    \item la evolución del producto total al aumentar un único factor variable.
    \item la recta isocosto asociada exclusivamente a los costos fijos medios.
    \item la función inversa de demanda enfrentada por la firma competitiva.
\end{opts}

\q{La demanda condicionada de un factor productivo se obtiene al:}
\begin{opts}
    \item maximizar el beneficio económico dado un precio de mercado.
    \item minimizar el costo para un nivel de producto y precios dados.
    \item maximizar la producción total sujeta a un presupuesto fijo.
    \item determinar simultáneamente la curva de oferta del mercado.
\end{opts}

\q{El lema de Shephard aplicado a la función de costos establece que:}
\begin{opts}
    \item $\partial C/\partial Q=CMg$.
    \item $\partial C/\partial w=L(w,r,Q)$.
    \item $\partial C/\partial L=w$.
    \item $\partial Q/\partial L=w$.
\end{opts}

\q{Existen economías de escala cuando, al aumentar la producción:}
\begin{opts}
    \item el costo medio aumenta.
    \item el costo marginal aumenta necesariamente.
    \item el costo medio disminuye.
    \item el costo total disminuye.
\end{opts}

\q{Existen economías de alcance cuando:}
\begin{opts}
    \item la empresa fabrica exclusivamente un único bien homogéneo.
    \item el costo medio de un bien cae al elevar su propia escala.
    \item el costo marginal resulta constante en todo nivel de $Q$.
    \item producir dos bienes juntos cuesta menos que por separado.
\end{opts}

\q{Un costo hundido es aquel que:}
\begin{opts}
    \item puede recuperarse íntegramente si la firma decide cerrar.
    \item coincide exactamente con el costo marginal de corto plazo.
    \item ya ocurrió, es irrecuperable y no debe guiar decisiones futuras.
    \item mide el costo de oportunidad del capital financiero propio.
\end{opts}

\q{El costo económico de producción se diferencia del costo contable porque:}
\begin{opts}
    \item incluye única y exclusivamente los egresos contables explícitos.
    \item excluye sistemáticamente la retribución normal al capital propio.
    \item incluye tanto los costos explícitos como los de oportunidad.
    \item excluye la depreciación económica de activos fijos productivos.
\end{opts}

\q{El beneficio normal se considera:}
\begin{opts}
    \item parte del costo económico al remunerar el capital propio.
    \item una ganancia extraordinaria superior al costo de oportunidad.
    \item un costo fijo contable irrelevante para la permanencia.
    \item una renta económica pura derivada de barreras de entrada.
\end{opts}

\q{En el corto plazo, una empresa competitiva decide suspender la producción cuando:}
\begin{opts}
    \item $P<CVMe$.
    \item $P<CMe$.
    \item $P<CMg$.
    \item $P<CFMe$.
\end{opts}

\q{En el largo plazo, una empresa competitiva decide salir del mercado cuando:}
\begin{opts}
    \item $P<CVMe$.
    \item $P<CMe$.
    \item $P<CMg$.
    \item $P<CFMe$.
\end{opts}

\q{La condición de maximización de beneficios de una empresa perfectamente competitiva es:}
\begin{opts}
    \item $P=CMe$.
    \item $IMg=CMe$.
    \item $P=CMg$.
    \item $P=CVMe$.
\end{opts}

\q{La curva de oferta de una empresa competitiva en el corto plazo coincide con:}
\begin{opts}
    \item toda la curva de costo medio total en su tramo decreciente.
    \item el tramo del costo marginal sobre el mínimo del $CVMe$.
    \item toda la curva de costo variable medio a partir del origen.
    \item la curva inversa de demanda enfrentada en el mercado local.
\end{opts}

\q{En el equilibrio de largo plazo de una industria perfectamente competitiva, se cumple:}
\begin{opts}
    \item $P=CMg=CMe$ con beneficio económico nulo.
    \item $P>CMe$ con beneficio económico positivo.
    \item $P<CVMe$ con cierre inmediato de planta.
    \item $P=CMg>CMe$ con rentas extraordinarias.
\end{opts}

\q{El excedente del productor se define como:}
\begin{opts}
    \item el ingreso total menos el costo total.
    \item el beneficio contable.
    \item el excedente del consumidor.
    \item el ingreso total menos el costo variable.
\end{opts}

\q{La dualidad entre producción y costos implica que:}
\begin{opts}
    \item la función de costos resulta independiente de la tecnología.
    \item la función de producción solo queda definida en corto plazo.
    \item dicha relación solo se cumple en tecnologías Cobb-Douglas.
    \item costos y producción son representaciones duales de la tecnología.
\end{opts}

\q{Si una función de producción es homogénea de grado uno, el teorema de Euler establece que:}
\begin{opts}
    \item $Q=PMg_K\cdot K+PMg_L\cdot L$.
    \item $Q=PMg_K\cdot K-PMg_L\cdot L$.
    \item $Q=PMe_K\cdot K+PMe_L\cdot L$.
    \item $Q=w\cdot L+r\cdot K\cdot P$.
\end{opts}

\q{La elasticidad de sustitución entre factores se define como:}
\begin{opts}
    \item el cambio porcentual del producto dividido por el cambio porcentual de los insumos.
    \item el cambio porcentual de la relación $K/L$ dividido por el cambio porcentual de la RMST.
    \item el cambio porcentual del costo dividido por el cambio porcentual del producto.
    \item el cambio porcentual del precio dividido por el cambio porcentual de la cantidad.
\end{opts}

\section*{Bloque B. Aplicaciones cuantitativas y optimización}

\noindent\textit{Preguntas 31 a 65. Requieren plantear y resolver problemas de producción, minimización de costos, maximización de beneficios y equilibrio competitivo con distintas formas funcionales, así como calcular productividades, elasticidades y economías de escala.}

\q{Con $Q=K^{0.5}L^{0.5}$, $K=16$ y $L=25$, la producción es:}
\begin{opts}
    \item 12
    \item 16
    \item 20
    \item 25
\end{opts}

\q{Con $Q=2K^{0.5}L^{0.5}$, $K=9$ y $L=4$, la producción es:}
\begin{opts}
    \item 12
    \item 24
    \item 6
    \item 18
\end{opts}

\q{Con $Q=K^{0.5}L^{0.5}$, $w=4$, $r=1$ y una meta de producción $Q=20$, el costo mínimo es:}
\begin{opts}
    \item 40
    \item 60
    \item 80
    \item 100
\end{opts}

\q{En el problema anterior, la relación óptima $K/L$ es:}
\begin{opts}
    \item 4
    \item $1/4$
    \item 2
    \item 1
\end{opts}

\q{Para $Q=K^{0.3}L^{0.6}$, la tecnología presenta rendimientos a escala:}
\begin{opts}
    \item crecientes (1.1).
    \item constantes (1.0).
    \item nulos (grado 0.0).
    \item decrecientes (0.9).
\end{opts}

\q{Para $Q=K^{0.4}L^{0.7}$, la tecnología presenta rendimientos a escala:}
\begin{opts}
    \item decrecientes.
    \item constantes.
    \item de 0.9.
    \item crecientes.
\end{opts}

\q{Con $Q=\min(2K,L)$, $w=2$, $r=3$ y $Q=40$, el costo mínimo es:}
\begin{opts}
    \item 100
    \item 120
    \item 140
    \item 160
\end{opts}

\q{Con $Q=\min(K,3L)$, $w=1$, $r=2$ y $Q=30$, el costo mínimo es:}
\begin{opts}
    \item 40
    \item 50
    \item 60
    \item 70
\end{opts}

\q{Con $Q=K+L$, $w=2$, $r=1$ y $Q=50$, el costo mínimo es:}
\begin{opts}
    \item 100
    \item 75
    \item 50
    \item 25
\end{opts}

\q{Con $Q=2K+L$, $w=3$, $r=2$ y $Q=40$, la cantidad óptima de capital es:}
\begin{opts}
    \item 10
    \item 15
    \item 20
    \item 40
\end{opts}

\q{En el corto plazo, $Q=10L^2-L^3$. El producto marginal del trabajo en $L=4$ es:}
\begin{opts}
    \item 32
    \item 24
    \item 16
    \item 8
\end{opts}

\q{Con la misma función $Q=10L^2-L^3$, el producto medio en $L=4$ es:}
\begin{opts}
    \item 16
    \item 24
    \item 32
    \item 20
\end{opts}

\q{Con $Q=10L^2-L^3$, el nivel de trabajo que maximiza el producto medio es:}
\begin{opts}
    \item 3
    \item 4
    \item 5
    \item 6
\end{opts}

\q{Si $C(Q)=100+4Q+Q^2$, el costo marginal en $Q=10$ es:}
\begin{opts}
    \item 4
    \item 14
    \item 20
    \item 24
\end{opts}

\q{Si $C(Q)=100+4Q+Q^2$, el costo medio en $Q=10$ es:}
\begin{opts}
    \item 14
    \item 20
    \item 22
    \item 24
\end{opts}

\q{Si $C(Q)=100+4Q+Q^2$, el costo variable medio en $Q=10$ es:}
\begin{opts}
    \item 4
    \item 14
    \item 24
    \item 10
\end{opts}

\q{Si $C(Q)=100+4Q+Q^2$, el nivel de producción que minimiza el costo medio es:}
\begin{opts}
    \item 5
    \item 20
    \item 14.14
    \item 10
\end{opts}

\q{Si $P=20$ y $C(Q)=100+4Q+Q^2$, la cantidad que maximiza el beneficio es:}
\begin{opts}
    \item 8
    \item 10
    \item 12
    \item 6
\end{opts}

\q{En el problema anterior, el beneficio económico en el óptimo es:}
\begin{opts}
    \item 36
    \item $-36$
    \item 0
    \item $-64$
\end{opts}

\q{Con $C(Q)=100+4Q+Q^2$, el precio de cierre o suspensión de la empresa en el corto plazo es:}
\begin{opts}
    \item 0
    \item 2
    \item 4
    \item 24
\end{opts}

\q{Suponga $C(Q)=100+4Q+Q^2$ para cada empresa y una demanda de mercado $Q=500-10P$. En el equilibrio de largo plazo, el número de empresas es:}
\begin{opts}
    \item 20
    \item 26
    \item 30
    \item 13
\end{opts}

\q{Si $P=20$, $CMg=4+2Q$ y la empresa produce $Q=8$, el excedente del productor es:}
\begin{opts}
    \item 32
    \item 64
    \item 96
    \item 128
\end{opts}

\q{Para $Q=(K^{0.5}+L^{0.5})^2$, la elasticidad de sustitución es:}
\begin{opts}
    \item 2
    \item 0.5
    \item 1
    \item infinita
\end{opts}

\q{Para $Q=(K^{-1}+L^{-1})^{-1}$, la elasticidad de sustitución es:}
\begin{opts}
    \item 0.0
    \item 1.0
    \item 2.0
    \item 0.5
\end{opts}

\q{Para $Q=K^{0.5}L^{0.5}$, con $w=4$ y $r=1$, la función de costo de largo plazo es:}
\begin{opts}
    \item $C=4Q$
    \item $C=2Q$
    \item $C=Q$
    \item $C=8Q$
\end{opts}

\q{Con $Q=K^{0.5}L^{0.5}$, $K=16$ fijo, $w=4$, $r=1$ y $Q=8$, el costo de corto plazo es:}
\begin{opts}
    \item 16
    \item 32
    \item 48
    \item 64
\end{opts}

\q{En el problema anterior, el costo marginal de corto plazo en $Q=8$ es:}
\begin{opts}
    \item 2
    \item 4
    \item 8
    \item 16
\end{opts}

\q{Si la función de costo de largo plazo es $C=4Q$, el costo marginal de largo plazo es:}
\begin{opts}
    \item 2
    \item 8
    \item 0
    \item 4
\end{opts}

\q{Suponga que el costo unitario es $C(Q)=100Q^{-0.2}$. Al duplicar la producción acumulada, el costo unitario se reduce aproximadamente en:}
\begin{opts}
    \item 5\%
    \item 10\%
    \item 13\%
    \item 20\%
\end{opts}

\q{Para $Q=K^{0.5}L^{0.5}$, la elasticidad del producto respecto del trabajo es:}
\begin{opts}
    \item 1
    \item 0.3
    \item 0.6
    \item 0.5
\end{opts}

\q{Si una tecnología tiene rendimientos crecientes a escala, la elasticidad del costo total respecto de la producción es:}
\begin{opts}
    \item mayor que 1.
    \item igual a 1.
    \item igual a 0.
    \item menor que 1.
\end{opts}

\q{Para $Q=K^{0.5}L^{0.5}$, la elasticidad de la demanda condicionada de trabajo respecto del precio del capital es:}
\begin{opts}
    \item $-0.5$.
    \item $0.5$.
    \item 1.
    \item $-1$.
\end{opts}

\q{Si $Q=AK^{0.5}L^{0.5}$ y la productividad total de los factores $A$ aumenta 10\%, manteniendo $K$ y $L$ constantes, la producción aumenta:}
\begin{opts}
    \item 5\%.
    \item 0\%.
    \item 10\%.
    \item 20\%.
\end{opts}

\q{Con $Q=K^{0.5}L^{0.5}$, $w=1$, $r=4$ y $Q=10$, el costo mínimo es:}
\begin{opts}
    \item 40
    \item 20
    \item 60
    \item 80
\end{opts}

\q{Con $Q=K^{0.5}L^{0.5}$, $K=16$ fijo, $w=4$, $r=1$ y $Q=20$, el costo de corto plazo es:}
\begin{opts}
    \item 100
    \item 116
    \item 125
    \item 80
\end{opts}

\section*{Bloque C. Casos con implicancias tarifarias y regulatorias del sector saneamiento}

\noindent\textit{Preguntas 66 a 90. Aplican la teoría del productor y de los costos a la interpretación de estructuras tarifarias, economías de escala y alcance, eficiencia productiva, productividad, agua no facturada y dilemas regulatorios de las empresas prestadoras de servicios de saneamiento (EPS).}

\q{Una EPS tiene una función de costo total $C(Q)=1000+2Q+0.1Q^2$, donde $Q$ se mide en miles de m$^3$. Si $Q=50$, el costo medio es:}
\begin{opts}
    \item 22
    \item 27
    \item 12
    \item 30
\end{opts}

\q{Con la misma función $C(Q)=1000+2Q+0.1Q^2$, el costo marginal en $Q=50$ es:}
\begin{opts}
    \item 2
    \item 7
    \item 22
    \item 12
\end{opts}

\q{Con $C(Q)=1000+2Q+0.1Q^2$, el nivel de producción que minimiza el costo medio es:}
\begin{opts}
    \item 50
    \item 70
    \item 100
    \item 150
\end{opts}

\q{Si el regulador fija una tarifa igual al costo marginal y la EPS produce $Q=50$, la pérdida que debe cubrirse con un subsidio o cargo fijo es:}
\begin{opts}
    \item 350
    \item 600
    \item 750
    \item 1350
\end{opts}

\q{Una empresa de saneamiento es un monopolio natural cuando:}
\begin{opts}
    \item el costo medio decrece y la función de costos es subaditiva.
    \item el costo medio es creciente en todo nivel de producción.
    \item el costo marginal es constante e igual al costo medio total.
    \item existen severas deseconomías de alcance entre los servicios.
\end{opts}

\q{Existen economías de alcance en una EPS que presta agua potable y alcantarillado cuando:}
\begin{opts}
    \item el costo medio del agua cae al ampliar solo su propio volumen.
    \item el costo marginal permanece constante en todo rango operativo.
    \item prestar ambos servicios juntos cuesta menos que por separado.
    \item los costos fijos de infraestructura resultan iguales a cero.
\end{opts}

\q{La regulación por precio tope RPI-X incentiva a la empresa a reducir costos porque:}
\begin{opts}
    \item retiene las ganancias de eficiencia dentro del período regulado.
    \item el regulador reembolsa automáticamente cualquier costo incurrido.
    \item puede sobrecapitalizar para ampliar la base tarifaria regulada.
    \item la tarifa anual se iguala siempre al costo medio observado.
\end{opts}

\q{El efecto Averch-Johnson, asociado a la regulación por tasa de retorno, tiende a generar:}
\begin{opts}
    \item sobrecapitalización para ampliar la base de capital regulada.
    \item subcapitalización extrema al no remunerar el capital invertido.
    \item competencia perfecta inmediata en el sector de saneamiento.
    \item eliminación total de los costos hundidos de infraestructura.
\end{opts}

\q{La regulación por costo de servicio (cost-plus) presenta el riesgo de:}
\begin{opts}
    \item incentivos sumamente fuertes para minimizar costos internos.
    \item tarifas fijas totalmente independientes del costo observado.
    \item competencia directa por comparación entre distintas firmas.
    \item incentivos débiles a la eficiencia y sesgo a inflar costos.
\end{opts}

\q{La competencia por comparación o yardstick competition consiste en:}
\begin{opts}
    \item fijar la tarifa uniforme exactamente igual al costo marginal.
    \item comparar el desempeño con empresas pares para fijar metas.
    \item regular exclusivamente la tasa de retorno sobre activos.
    \item permitir a la EPS cobrar libremente el precio de monopolio.
\end{opts}

\q{Si el precio de la energía eléctrica, insumo intensivo en el bombeo de agua, aumenta de forma permanente, en el largo plazo la EPS tenderá a:}
\begin{opts}
    \item sustituir energía por capital mediante equipos eficientes.
    \item sustituir capital físico por mayor uso intensivo de energía.
    \item sustituir mano de obra calificada por energía eléctrica.
    \item mantener rígida la misma proporción histórica de insumos.
\end{opts}

\q{Una EPS produce 1,000,000 de m$^3$ al año con 50 trabajadores. Su productividad media del trabajo es:}
\begin{opts}
    \item 10,000 m$^3$ por trabajador.
    \item 50,000 m$^3$ por trabajador.
    \item 5,000 m$^3$ por trabajador.
    \item 20,000 m$^3$ por trabajador.
\end{opts}

\q{Si el producto de una EPS crece 6\% y sus insumos crecen 4\%, la productividad total de los factores (residual de Solow) crece aproximadamente:}
\begin{opts}
    \item 10\%.
    \item 6\%.
    \item 2\%.
    \item 4\%.
\end{opts}

\q{Si el agua no facturada (ANF) es 25\% del agua producida y el costo de producción es S/ 1.50 por m$^3$, el costo efectivo por m$^3$ facturado es:}
\begin{opts}
    \item S/ 1.50.
    \item S/ 1.875.
    \item S/ 2.25.
    \item S/ 2.00.
\end{opts}

\q{Si la ANF se reduce de 25\% a 15\%, manteniendo el costo de producción en S/ 1.50 por m$^3$, el costo efectivo por m$^3$ facturado se reduce en aproximadamente:}
\begin{opts}
    \item S/ 0.12.
    \item S/ 0.24.
    \item S/ 0.35.
    \item S/ 0.50.
\end{opts}

\q{Un programa de reducción de fugas cuesta S/ 80,000 al año y permite ahorrar 50,000 m$^3$ de agua cuyo costo marginal de producción es S/ 2 por m$^3$. El beneficio neto anual del programa es:}
\begin{opts}
    \item S/ 0.
    \item S/ 10,000.
    \item S/ 20,000.
    \item S/ 30,000.
\end{opts}

\q{En una tarifa por bloques horarios o de punta, la lógica de eficiencia indica que:}
\begin{opts}
    \item se cobre una tarifa idéntica en todas las franjas horarias.
    \item se cobre más en hora punta por el mayor costo marginal de capacidad.
    \item se cobre menos en hora punta para estimular el consumo nocturno.
    \item se fije tarifa cero durante los períodos de máxima congestión.
\end{opts}

\q{Cuando una inversión en redes es hundida y específica, el regulador puede enfrentar un problema de:}
\begin{opts}
    \item compromiso creíble: riesgo de no reconocimiento tras invertir.
    \item competencia perfecta entre múltiples redes superpuestas.
    \item ausencia completa de costos fijos en la etapa de distribución.
    \item inexistencia de economías de escala en tuberías principales.
\end{opts}

\q{Si la base de activos regulados es S/ 1,000,000 y la depreciación lineal es 5\% anual, la depreciación anual es:}
\begin{opts}
    \item S/ 50,000.
    \item S/ 100,000.
    \item S/ 5,000.
    \item S/ 20,000.
\end{opts}

\q{Un costo común de S/ 100,000 se reparte entre agua y alcantarillado según sus costos separados: 60\% agua y 40\% alcantarillado. La parte asignada al agua es:}
\begin{opts}
    \item S/ 40,000.
    \item S/ 50,000.
    \item S/ 60,000.
    \item S/ 100,000.
\end{opts}

\q{Si el costo fijo anual de una EPS es S/ 1,000,000 y tiene 10,000 usuarios, el cargo fijo puro por usuario para cubrirlo es:}
\begin{opts}
    \item S/ 10.
    \item S/ 50.
    \item S/ 1,000.
    \item S/ 100.
\end{opts}

\q{Si el costo fijo que debe recuperarse es S/ 500,000 y el volumen facturado es 100,000 m$^3$, el cargo volumétrico uniforme que lo cubre es:}
\begin{opts}
    \item S/ 0.50 por m$^3$.
    \item S/ 1.00 por m$^3$.
    \item S/ 2.00 por m$^3$.
    \item S/ 5.00 por m$^3$.
\end{opts}

\q{Si dos EPS pequeñas se fusionan y existen economías de escala en el rango relevante, cabe esperar que la fusión:}
\begin{opts}
    \item baje el costo medio si las economías superan la coordinación.
    \item suba inevitablemente el costo medio desde la primera etapa.
    \item carezca de todo impacto sobre los costos medios unitarios.
    \item solo funcione cuando existan economías de alcance y no escala.
\end{opts}

\q{Si el regulador exige mejorar la continuidad y la presión del servicio, la teoría de costos sugiere que:}
\begin{opts}
    \item mayor calidad eleva costos y exige reconocimiento tarifario.
    \item la mejora técnica es gratuita y no altera los costos reales.
    \item reducir calidad siempre baja costos sin dañar el bienestar.
    \item la calidad operativa es ajena a la inversión y mantenimiento.
\end{opts}

\q{Si el costo marginal de conservar agua mediante reducción de fugas es S/ 1.50 por m$^3$ y el costo marginal de ampliar una nueva fuente es S/ 2.50 por m$^3$, la opción económicamente preferible es:}
\begin{opts}
    \item ampliar de inmediato la nueva fuente.
    \item ambas alternativas tienen costo igual.
    \item mantener las fugas sin intervención.
    \item priorizar la reducción de pérdidas.
\end{opts}

\newpage
\section*{Tabla de claves rápidas}

\noindent Use esta tabla para una verificación inmediata (preguntas 1 a 180). El solucionario con la justificación de cada respuesta se presenta en la sección siguiente.

\vspace{0.35cm}
\begin{center}
\footnotesize
\renewcommand{\arraystretch}{1.25}
\begin{tabular}{|c|c|c|c|c|c|c|c|c|c|c|}
\hline
\textbf{N.} & \textbf{1} & \textbf{2} & \textbf{3} & \textbf{4} & \textbf{5} & \textbf{6} & \textbf{7} & \textbf{8} & \textbf{9} & \textbf{10} \\
\hline
\textbf{Clave} & B & C & A & C & C & D & B & B & D & C \\
\hline
\textbf{N.} & \textbf{11} & \textbf{12} & \textbf{13} & \textbf{14} & \textbf{15} & \textbf{16} & \textbf{17} & \textbf{18} & \textbf{19} & \textbf{20} \\
\hline
\textbf{Clave} & B & B & A & A & B & B & C & D & C & C \\
\hline
\textbf{N.} & \textbf{21} & \textbf{22} & \textbf{23} & \textbf{24} & \textbf{25} & \textbf{26} & \textbf{27} & \textbf{28} & \textbf{29} & \textbf{30} \\
\hline
\textbf{Clave} & A & A & B & C & B & A & D & D & A & B \\
\hline
\textbf{N.} & \textbf{31} & \textbf{32} & \textbf{33} & \textbf{34} & \textbf{35} & \textbf{36} & \textbf{37} & \textbf{38} & \textbf{39} & \textbf{40} \\
\hline
\textbf{Clave} & C & A & C & A & D & D & C & D & C & C \\
\hline
\textbf{N.} & \textbf{41} & \textbf{42} & \textbf{43} & \textbf{44} & \textbf{45} & \textbf{46} & \textbf{47} & \textbf{48} & \textbf{49} & \textbf{50} \\
\hline
\textbf{Clave} & A & B & C & D & D & B & D & A & B & C \\
\hline
\textbf{N.} & \textbf{51} & \textbf{52} & \textbf{53} & \textbf{54} & \textbf{55} & \textbf{56} & \textbf{57} & \textbf{58} & \textbf{59} & \textbf{60} \\
\hline
\textbf{Clave} & B & B & A & D & A & B & B & D & C & D \\
\hline
\textbf{N.} & \textbf{61} & \textbf{62} & \textbf{63} & \textbf{64} & \textbf{65} & \textbf{66} & \textbf{67} & \textbf{68} & \textbf{69} & \textbf{70} \\
\hline
\textbf{Clave} & D & B & C & A & B & B & D & C & C & A \\
\hline
\textbf{N.} & \textbf{71} & \textbf{72} & \textbf{73} & \textbf{74} & \textbf{75} & \textbf{76} & \textbf{77} & \textbf{78} & \textbf{79} & \textbf{80} \\
\hline
\textbf{Clave} & C & A & A & D & B & A & D & C & D & B \\
\hline
\textbf{N.} & \textbf{81} & \textbf{82} & \textbf{83} & \textbf{84} & \textbf{85} & \textbf{86} & \textbf{87} & \textbf{88} & \textbf{89} & \textbf{90} \\
\hline
\textbf{Clave} & C & B & A & A & C & D & D & A & A & D \\
\hline
\textbf{N.} & \textbf{91} & \textbf{92} & \textbf{93} & \textbf{94} & \textbf{95} & \textbf{96} & \textbf{97} & \textbf{98} & \textbf{99} & \textbf{100} \\
\hline
\textbf{Clave} & C & A & D & B & D & C & A & B & C & D \\
\hline
\textbf{N.} & \textbf{101} & \textbf{102} & \textbf{103} & \textbf{104} & \textbf{105} & \textbf{106} & \textbf{107} & \textbf{108} & \textbf{109} & \textbf{110} \\
\hline
\textbf{Clave} & B & A & C & B & D & A & B & C & A & D \\
\hline
\textbf{N.} & \textbf{111} & \textbf{112} & \textbf{113} & \textbf{114} & \textbf{115} & \textbf{116} & \textbf{117} & \textbf{118} & \textbf{119} & \textbf{120} \\
\hline
\textbf{Clave} & C & B & A & D & B & C & D & A & B & C \\
\hline
\textbf{N.} & \textbf{121} & \textbf{122} & \textbf{123} & \textbf{124} & \textbf{125} & \textbf{126} & \textbf{127} & \textbf{128} & \textbf{129} & \textbf{130} \\
\hline
\textbf{Clave} & D & A & B & C & A & D & C & B & A & D \\
\hline
\textbf{N.} & \textbf{131} & \textbf{132} & \textbf{133} & \textbf{134} & \textbf{135} & \textbf{136} & \textbf{137} & \textbf{138} & \textbf{139} & \textbf{140} \\
\hline
\textbf{Clave} & B & C & D & A & C & B & D & A & B & C \\
\hline
\textbf{N.} & \textbf{141} & \textbf{142} & \textbf{143} & \textbf{144} & \textbf{145} & \textbf{146} & \textbf{147} & \textbf{148} & \textbf{149} & \textbf{150} \\
\hline
\textbf{Clave} & A & D & B & C & D & A & C & B & A & D \\
\hline
\textbf{N.} & \textbf{151} & \textbf{152} & \textbf{153} & \textbf{154} & \textbf{155} & \textbf{156} & \textbf{157} & \textbf{158} & \textbf{159} & \textbf{160} \\
\hline
\textbf{Clave} & C & B & D & A & B & C & A & D & C & B \\
\hline
\textbf{N.} & \textbf{161} & \textbf{162} & \textbf{163} & \textbf{164} & \textbf{165} & \textbf{166} & \textbf{167} & \textbf{168} & \textbf{169} & \textbf{170} \\
\hline
\textbf{Clave} & A & D & B & C & A & D & B & C & A & B \\
\hline
\textbf{N.} & \textbf{171} & \textbf{172} & \textbf{173} & \textbf{174} & \textbf{175} & \textbf{176} & \textbf{177} & \textbf{178} & \textbf{179} & \textbf{180} \\
\hline
\textbf{Clave} & D & C & B & A & C & D & B & C & D & A \\
\hline
\end{tabular}
\end{center}

\newpage
\section*{Solucionario completo}

\noindent Cada respuesta incluye la clave y una justificación breve. En los ejercicios numéricos se indica la ecuación o el cálculo principal.

\subsection*{Bloque A. Fundamentos teóricos y conceptuales}

\sol{1}{B}{La función de producción describe la frontera tecnológica: la cantidad máxima de producto obtenible con combinaciones dadas de insumos. No es una función de costos ni de beneficios.}

\sol{2}{C}{Una isocuanta agrupa todas las combinaciones de $K$ y $L$ que generan el mismo nivel de producción. La isocosto, en cambio, agrupa combinaciones de igual costo.}

\sol{3}{A}{La RMST indica cuánto capital puede reducirse al aumentar una unidad de trabajo sin alterar la producción. Es la pendiente (en valor absoluto) de la isocuanta.}

\sol{4}{C}{$RMST_{LK}=PMg_L/PMg_K$. En el óptimo de costos se iguala a $w/r$.}

\sol{5}{C}{Una RMST decreciente significa que cada vez se necesita menos capital para compensar una unidad adicional de trabajo; las isocuantas son convexas al origen.}

\sol{6}{D}{Rendimientos constantes a escala: $f(\lambda K,\lambda L)=\lambda f(K,L)$. Al duplicar los insumos, el producto se duplica.}

\sol{7}{B}{En $Q=AK^\alpha L^\beta$, al multiplicar los insumos por $\lambda$ se obtiene $\lambda^{\alpha+\beta}Q$. El grado de homogeneidad es $\alpha+\beta$.}

\sol{8}{B}{Si $\alpha+\beta>1$, el producto crece más que proporcionalmente al aumentar todos los insumos: rendimientos crecientes a escala.}

\sol{9}{D}{La ley de rendimientos marginales decrecientes es un fenómeno de corto plazo: con al menos un factor fijo, el producto marginal del factor variable eventualmente disminuye.}

\sol{10}{C}{El producto medio alcanza su máximo cuando se iguala al producto marginal. Si $PMg>PMe$, el $PMe$ aumenta; si $PMg<PMe$, disminuye.}

\sol{11}{B}{La etapa II comienza en el máximo del producto medio y termina cuando el producto marginal se hace cero. Es la región económicamente relevante.}

\sol{12}{B}{De $C=wL+rK$ se obtiene $K=C/r-(w/r)L$; la pendiente es $-w/r$.}

\sol{13}{A}{En el óptimo interior, la isocuanta es tangente a la isocosto: $RMST_{LK}=w/r$.}

\sol{14}{A}{La senda de expansión une las combinaciones óptimas de insumos para distintos niveles de producción, manteniendo los precios de los insumos constantes.}

\sol{15}{B}{La demanda condicionada de factores se deriva de minimizar $wL+rK$ sujeto a $Q=f(K,L)$. Depende de $Q$, $w$ y $r$.}

\sol{16}{B}{El lema de Shephard para costos: $\partial C(w,r,Q)/\partial w=L(w,r,Q)$. La derivada del costo respecto del precio del insumo es la demanda condicionada.}

\sol{17}{C}{Economías de escala: el costo medio disminuye al aumentar la producción. No implica que el costo total disminuya.}

\sol{18}{D}{Economías de alcance: $C(Q_1,Q_2)<C(Q_1,0)+C(0,Q_2)$. Producir conjuntamente es más barato que hacerlo por separado.}

\sol{19}{C}{Un costo hundido ya se incurrió y no es recuperable. No debe influir en decisiones marginales futuras, aunque puede afectar la decisión inicial de entrar.}

\sol{20}{C}{El costo económico incluye costos explícitos (pagos efectivos) y costos implícitos o de oportunidad, como el uso de capital propio o el tiempo del empresario.}

\sol{21}{A}{El beneficio normal es el rendimiento mínimo exigido por el capital propio; forma parte del costo de oportunidad y, por ello, del costo económico.}

\sol{22}{A}{En el corto plazo, si $P<CVMe$, la empresa no cubre ni los costos variables y le conviene suspender. Si $CVMe<P<CMe$, cubre variables y parte de los fijos, por lo que sigue operando.}

\sol{23}{B}{En el largo plazo no hay costos fijos irreversibles. Si $P<CMe$, la empresa obtiene pérdidas y sale del mercado.}

\sol{24}{C}{La empresa competitiva enfrenta $P=IMg$; maximiza beneficio donde $IMg=CMg$, es decir, $P=CMg$.}

\sol{25}{B}{La curva de oferta de corto plazo es el tramo de $CMg$ por encima del mínimo $CVMe$. Por debajo de ese precio, la empresa suspende.}

\sol{26}{A}{En el equilibrio de largo plazo, la libre entrada y salida lleva a $P=CMg=CMe$ y beneficio económico nulo; las empresas obtienen solo beneficio normal.}

\sol{27}{D}{El excedente del productor es el ingreso total menos el costo variable. Equivale al beneficio económico más el costo fijo.}

\sol{28}{D}{La dualidad establece que la función de costos contiene la misma información tecnológica que la función de producción, bajo condiciones de regularidad.}

\sol{29}{A}{Si $Q$ es homogénea de grado uno, Euler implica $Q=PMg_K K+PMg_L L$: el producto se agota si cada factor recibe su producto marginal.}

\sol{30}{B}{La elasticidad de sustitución mide la respuesta porcentual de la relación $K/L$ ante un cambio porcentual de la RMST: $\sigma=\dfrac{\%\Delta(K/L)}{\%\Delta RMST}$.}

\subsection*{Bloque B. Aplicaciones cuantitativas y optimización}

\sol{31}{C}{$Q=16^{0.5}25^{0.5}=4\cdot5=20$.}

\sol{32}{A}{$Q=2(9^{0.5})(4^{0.5})=2(3)(2)=12$.}

\sol{33}{C}{En el óptimo, $RMST=PMg_L/PMg_K=K/L=w/r=4$, de donde $K=4L$. Con $Q=20$: $Q=\sqrt{4L^2}=2L=20$, entonces $L=10$ y $K=40$. El costo es $4(10)+1(40)=80$.}

\sol{34}{A}{$K/L=w/r=4/1=4$.}

\sol{35}{D}{La suma de exponentes es $0.3+0.6=0.9<1$: rendimientos decrecientes a escala.}

\sol{36}{D}{La suma de exponentes es $0.4+0.7=1.1>1$: rendimientos crecientes a escala.}

\sol{37}{C}{En Leontief, $2K=L=40$. Entonces $K=20$ y $L=40$. El costo es $2(40)+3(20)=80+60=140$.}

\sol{38}{D}{Se requiere $K=30$ y $3L=30$, es decir $L=10$. El costo es $1(10)+2(30)=10+60=70$.}

\sol{39}{C}{El capital es más barato por unidad ($r=1<w=2$). Se usa solo capital: $K=50$, $L=0$; costo $=50$.}

\sol{40}{C}{$PMg_K/r=2/2=1$ y $PMg_L/w=1/3$. Conviene usar solo capital: $2K=40$, entonces $K=20$ y el costo es $2(20)=40$.}

\sol{41}{A}{$PMg_L=20L-3L^2$. En $L=4$: $80-48=32$.}

\sol{42}{B}{$Q=10(16)-64=96$; $PMe=Q/L=96/4=24$.}

\sol{43}{C}{$PMe=Q/L=10L-L^2$. Su máximo se obtiene con $dPMe/dL=10-2L=0$, es decir $L=5$.}

\sol{44}{D}{$CMg=dC/dQ=4+2Q$. En $Q=10$: $4+20=24$.}

\sol{45}{D}{$CMe=C/Q=100/Q+4+Q$. En $Q=10$: $10+4+10=24$.}

\sol{46}{B}{$CV=4Q+Q^2$. En $Q=10$: $40+100=140$; $CVMe=140/10=14$.}

\sol{47}{D}{$CMe=100/Q+4+Q$. Derivando: $-100/Q^2+1=0$, de donde $Q^2=100$ y $Q=10$.}

\sol{48}{A}{$P=CMg$: $20=4+2Q$, entonces $Q=8$.}

\sol{49}{B}{$IT=20(8)=160$. $CT=100+4(8)+8^2=100+32+64=196$. Beneficio $=160-196=-36$.}

\sol{50}{C}{$CVMe=4+Q$, que se minimiza en $Q=0$ con valor 4. Si $P<4$, la empresa no cubre ni los costos variables y suspende.}

\sol{51}{B}{En largo plazo, $P=\min CMe=24$ y cada empresa produce $Q=10$. La demanda de mercado es $Q=500-10(24)=260$. El número de empresas es $260/10=26$.}

\sol{52}{B}{$PS=\int_0^8[20-(4+2Q)]dQ=16(8)-8^2=128-64=64$.}

\sol{53}{A}{La función es CES con $\rho=0.5$. Entonces $\sigma=1/(1-\rho)=1/0.5=2$.}

\sol{54}{D}{Aquí $\rho=-1$; $\sigma=1/(1-(-1))=0.5$.}

\sol{55}{A}{Para $Q=K^{0.5}L^{0.5}$, $C=2Q\sqrt{wr}$. Con $w=4$ y $r=1$: $C=2Q(2)=4Q$.}

\sol{56}{B}{Con $K=16$, $Q=4\sqrt{L}$. Para $Q=8$: $\sqrt{L}=2$, $L=4$. El costo es $4(4)+1(16)=32$.}

\sol{57}{B}{De $C=Q^2/4+16$, $CMg=Q/2$. En $Q=8$: $CMg=4$.}

\sol{58}{D}{Si $C=4Q$, el costo marginal es constante e igual a 4.}

\sol{59}{C}{El costo unitario satisface $C(2Q)/C(Q)=2^{-0.2}=0.8706$. Cae aproximadamente $1-0.8706=0.1294$, es decir, 13\%.}

\sol{60}{D}{En $Q=K^{0.5}L^{0.5}$, la elasticidad del producto respecto de $L$ es el exponente 0.5.}

\sol{61}{D}{Con rendimientos crecientes a escala, el costo total crece menos que proporcionalmente que la producción: la elasticidad costo-producto es menor que 1.}

\sol{62}{B}{$L=Q\sqrt{r/w}$. La elasticidad respecto de $r$ es $0.5$: un aumento de 1\% en $r$ eleva $L$ en 0.5\%.}

\sol{63}{C}{Si $A$ aumenta 10\% y $K,L$ no cambian, $Q$ aumenta exactamente 10\%, porque $Q$ es proporcional a $A$.}

\sol{64}{A}{$K/L=w/r=1/4$, entonces $K=0.25L$. $Q=\sqrt{0.25L^2}=0.5L=10$, de donde $L=20$ y $K=5$. El costo es $1(20)+4(5)=40$.}

\sol{65}{B}{Con $K=16$, $Q=4\sqrt{L}=20$, entonces $\sqrt{L}=5$ y $L=25$. El costo es $4(25)+1(16)=116$.}

\subsection*{Bloque C. Casos con implicancias tarifarias y regulatorias}

\sol{66}{B}{$C(50)=1000+2(50)+0.1(50)^2=1000+100+250=1350$. $CMe=1350/50=27$.}

\sol{67}{D}{$CMg=dC/dQ=2+0.2Q$. En $Q=50$: $2+10=12$.}

\sol{68}{C}{$CMe=1000/Q+2+0.1Q$. Derivando: $-1000/Q^2+0.1=0$, de donde $Q^2=10000$ y $Q=100$.}

\sol{69}{C}{Con $P=CMg=12$ y $Q=50$, el ingreso es $12(50)=600$ y el costo es $1350$. La pérdida es $1350-600=750$.}

\sol{70}{A}{Un monopolio natural se caracteriza por costos medios decrecientes en el rango relevante y una función de costos subaditiva: una sola empresa abastece el mercado a menor costo que varias.}

\sol{71}{C}{Las economías de alcance surgen cuando producir agua y alcantarillado conjuntamente es más barato que producir cada servicio por separado.}

\sol{72}{A}{En RPI-X, la tarifa se ajusta por inflación menos un factor de eficiencia. Si la empresa reduce costos por debajo de lo previsto, retiene la diferencia hasta el siguiente ajuste.}

\sol{73}{A}{El efecto Averch-Johnson sostiene que, al remunerar el capital, la empresa regulada tiende a sobrecapitalizar para elevar la base sobre la que obtiene retorno.}

\sol{74}{D}{En cost-plus, el regulador reconoce los costos incurridos; la empresa tiene pocos incentivos a reducirlos y puede inflarlos.}

\sol{75}{B}{El yardstick competition compara la empresa con un grupo de referencia o con empresas similares, introduciendo presión competitiva sin duplicar la red.}

\sol{76}{A}{Si la energía se encarece, en el largo plazo la empresa sustituye energía por capital (equipos más eficientes) para reducir el costo por m$^3$.}

\sol{77}{D}{Productividad media del trabajo $=Q/L=1,000,000/50=20,000$ m$^3$ por trabajador.}

\sol{78}{C}{Crecimiento de la PTF $\approx$ crecimiento del producto menos crecimiento de los insumos $=6\%-4\%=2\%$.}

\sol{79}{D}{Si 25\% del agua producida no se factura, se factura 75\%. El costo por m$^3$ facturado es $1.50/0.75=2.00$.}

\sol{80}{B}{Con ANF de 15\%, el costo por m$^3$ facturado es $1.50/0.85=1.7647$. La reducción es $2.00-1.7647=0.2353\approx0.24$.}

\sol{81}{C}{El ahorro valorizado es $50,000\times2=100,000$. El beneficio neto es $100,000-80,000=20,000$.}

\sol{82}{B}{En hora punta se requiere capacidad adicional y el costo marginal de suministro es mayor; la tarifa de punta debe reflejarlo para incentivar el uso eficiente.}

\sol{83}{A}{La inversión hundida crea un problema de compromiso: una vez realizada, la empresa no puede retirar el capital y queda expuesta a que el regulador no reconozca su costo.}

\sol{84}{A}{Depreciación $=0.05\times1,000,000=50,000$.}

\sol{85}{C}{La parte del costo común asignada al agua es $0.60\times100,000=60,000$.}

\sol{86}{D}{Cargo fijo por usuario $=1,000,000/10,000=100$.}

\sol{87}{D}{Cargo volumétrico $=500,000/100,000=5$ por m$^3$.}

\sol{88}{A}{La fusión puede reducir el costo medio si las economías de escala son suficientes para compensar los costos de coordinación y gestión.}

\sol{89}{A}{Mejorar continuidad y presión exige inversión, mantenimiento y operación; la regulación debe reconocer esos costos para preservar la sostenibilidad.}

\sol{90}{D}{La conservación cuesta S/ 1.50 por m$^3$, menos que los S/ 2.50 de la nueva fuente; es la opción económicamente preferible antes de ampliar la oferta.}

\subsection*{Bloque D. Ampliación: Teoría y Aplicaciones Prácticas (91 a 180)}

\sol{91}{C}{La libre disposición implica que si $y \in Y$ y $y' \le y$, entonces $y' \in Y$: la firma siempre puede descartar excedentes de producto o utilizar más insumos para obtener la misma o menor producción.}
\sol{92}{A}{Si $0 \in Y$ e $y \in Y$, por convexidad $\lambda y + (1-\lambda)0 = \lambda y \in Y$ para todo $\lambda \in [0,1]$, lo cual define rendimientos no crecientes a escala (se puede reducir la escala sin pérdida de eficiencia).}
\sol{93}{D}{Al mantener $\bar{K}$ fijo, cada unidad adicional de $L$ dispone de menos capital por trabajador, por lo que desde cierto punto $\partial^2 f/\partial L^2 < 0$ ($PMg_L$ decreciente). Esto es un fenómeno de corto plazo distinto a los rendimientos a escala.}
\sol{94}{B}{La etapa I termina donde $PMe_L$ es máximo ($PMg_L = PMe_L$) y la etapa III comienza cuando $PMg_L < 0$. En la etapa II, $PMe_L$ cae y $0 \le PMg_L \le PMe_L$ (elasticidad $\varepsilon_{Q,L} \in [0,1]$).}
\sol{95}{D}{Por definición, $\varepsilon_{Q,L} = \frac{\partial Q}{\partial L}\frac{L}{Q} = \frac{PMg_L}{PMe_L}$. En el máximo de $PMe_L$ se cumple $PMg_L = PMe_L$, por lo que $\varepsilon_{Q,L} = 1$.}
\sol{96}{C}{Por el teorema de Euler con homogeneidad de grado 1, $Q = L\cdot PMg_L + K\cdot PMg_K \implies PMe_L = PMg_L + \frac{K}{L}PMg_K$. Si $PMg_L > PMe_L$, necesariamente $PMg_K < 0$.}
\sol{97}{A}{Derivando $PMe_L = Q/L$ respecto a $L$ se obtiene $\frac{d PMe_L}{d L} = \frac{1}{L}(PMg_L - PMe_L)$, cuyo signo es estrictamente negativo whenever $PMg_L < PMe_L$.}
\sol{98}{B}{La convexidad estricta de las isocuantas (cuasiconcavidad estricta de $f$) implica que, conforme aumenta $L$ y cae $K$ sobre la misma isocuanta, la pendiente en valor absoluto $TMST_{L,K} = PMg_L/PMg_K$ decrece.}
\sol{99}{C}{Por definición, $\sigma = \frac{d\ln(K/L)}{d\ln(TMST_{L,K})}$, midiendo la curvatura de la isocuanta y cuán fácilmente la firma sustituye factores cuando cambian sus precios relativos $w/r$.}
\sol{100}{D}{En Cobb-Douglas, $TMST_{L,K} = \frac{\beta}{\alpha}\frac{K}{L} \implies \ln(K/L) = \ln(\alpha/\beta) + \ln(TMST_{L,K})$, por lo que $\sigma = \frac{d\ln(K/L)}{d\ln(TMST_{L,K})} = 1$ independientemente del grado de homogeneidad $\alpha+\beta$.}
\sol{101}{B}{Al multiplicar $(K,L)$ por $t>0$, $Q(tK, tL) = t^\nu Q(K,L)$. Así, $\nu$ gobierna los rendimientos a escala ($\nu=1$ constantes, $\nu>1$ crecientes), mientras que $\rho$ determina $\sigma = 1/(1-\rho)$.}
\sol{102}{A}{Como los insumos son complementos perfectos ($\sigma = 0$), el costo mínimo siempre ocurre en el vértice de cada isocuanta ($aK = bL$), de modo que la razón óptima $K/L = b/a$ es invariante a $w/r$.}
\sol{103}{C}{Con isocuantas lineales de pendiente constante $TMST_{L,K} = b/a$, si $b/a > w/r$ solo se usa $L$; si $b/a < w/r$ solo se usa $K$. Solo cuando $b/a = w/r$ coinciden isocuanta e isocosto permitiendo combinaciones interiores.}
\sol{104}{B}{La Translog es una aproximación de Taylor de segundo orden en logaritmos que no restringe $\sigma = 1$ ni los rendimientos a escala a ser constantes, permitiendo capturar economías de escala y alcance variables.}
\sol{105}{D}{Por el teorema de Euler, $K\cdot PMg_K + L\cdot PMg_L = k Q$. Multiplicando por $p$, el gasto en factores $rK + wL = k(pQ) > pQ$ cuando $k > 1$, lo que implica pérdidas si se fija $p = CMg$ (problema clásico del monopolio natural).}
\sol{106}{A}{En $Q = K^{0.7}L^{0.7}$, $\partial^2 Q/\partial L^2 < 0$ y $\partial^2 Q/\partial K^2 < 0$ (rendimientos marginales decrecientes en corto plazo), pero $Q(tK, tL) = t^{1.4}Q(K,L)$ (rendimientos crecientes a escala en largo plazo).}
\sol{107}{B}{Evaluando $\frac{d\ln f(tK, tL)}{d\ln t}\big|_{t=1} = \frac{K f_K + L f_L}{Q} = \varepsilon_{Q,K} + \varepsilon_{Q,L}$, se obtiene la elasticidad de escala local.}
\sol{108}{C}{La transformación monótona $F$ no altera la $TMST_{L,K}$ (que solo depende de $K/L$ por ser $g$ homogénea), por lo que la razón óptima $K^*/L^*$ es constante a lo largo de la senda de expansión, aunque $F$ permita rendimientos a escala variables.}
\sol{109}{A}{Con rendimientos constantes a escala, $C(w,r,q) = c(w,r)q$. Si $p > c(w,r)$, $\pi = (p-c)q \to \infty$ (no hay máximo finito); si $p = c(w,r)$, $\pi = 0$ para todo $q \ge 0$; si $p < c(w,r)$, $q^* = 0$.}
\sol{110}{D}{Por una propiedad estándar de funciones homogéneas, si $f$ es homogénea de grado $k$, sus derivadas de primer orden son homogéneas de grado $k-1$. Si $k=1$, $PMg_K$ y $PMg_L$ son homogéneas de grado $0$ (solo dependen de $K/L$).}
\sol{111}{C}{Reordenando $\frac{PMg_L}{PMg_K} = \frac{w}{r}$ se obtiene $\frac{PMg_L}{w} = \frac{PMg_K}{r}$, es decir, la productividad marginal por sol gastado se iguala en todos los factores utilizados.}
\sol{112}{B}{Por el teorema de la envolvente, $\frac{\partial C(w,r,q)}{\partial q} = \frac{\partial \mathcal{L}}{\partial q} = \mu^*$, interpretándose $\mu^*$ como el costo marginal de aumentar la meta de producción en una unidad.}
\sol{113}{A}{Si la firma mantuviera fijos los insumos iniciales $(L^0, K^0)$, el costo crecería linealmente con $(w,r)$; al reoptimizar y sustituir hacia el factor más barato, $C(w,r,q) \le wL^0 + rK^0$, lo que implica concavidad en $(w,r)$.}
\sol{114}{D}{Duplicar todos los precios de insumos no cambia la pendiente de las isocostos ($w/r$) ni la isocuanta objetivo $q$, por lo que la cesta óptima de insumos no varía (grado $0$) y el costo total se duplica (grado $1$).}
\sol{115}{B}{En el largo plazo la firma minimiza costos sobre un conjunto más amplio (sin la restricción $K=\bar{K}$), por lo que $C_{LP} \le C_{CP}$; ambos coinciden exactamente en el nivel de $q$ para el cual $\bar{K}$ es el capital óptimo de largo plazo.}
\sol{116}{C}{Si escalar insumos por $t$ multiplica el producto por $t^k = q$, entonces se requiere escalar los insumos (y el costo) por $t = q^{1/k}$, resultando $C(w,r,q) = C(w,r,1) \cdot q^{1/k}$.}
\sol{117}{D}{Por el lema de Shephard, $L^c = \partial C/\partial w$, de modo que $\partial L^c/\partial w = \partial^2 C/\partial w^2 \le 0$ por concavidad de $C$ en precios (ley de la demanda condicional de factores).}
\sol{118}{A}{Para minimizar una función lineal ($wL+rK$) sujeta a $f(K,L)=q$, basta que el conjunto superior de la isocuanta sea convexo (cuasiconcavidad de $f$), sin requerir concavidad global de $f$.}
\sol{119}{B}{Por el teorema de la envolvente aplicado al problema de minimización de costos, $\frac{\partial C(w,r,q)}{\partial w} = L^c(w,r,q)$ y $\frac{\partial C(w,r,q)}{\partial r} = K^c(w,r,q)$.}
\sol{120}{C}{Dado que $L^c = \partial C/\partial w$ y $K^c = \partial C/\partial r$, las derivadas cruzadas coinciden con la segunda derivada mixta $\frac{\partial^2 C}{\partial w \partial r} = \frac{\partial^2 C}{\partial r \partial w}$.}
\sol{121}{D}{Multiplicar $(p,w,r)$ por $t>0$ no altera los planes óptimos y multiplica $\pi$ por $t$ (grado 1). Además, $\pi(p,w,r)$ es el supremo de funciones lineales en $(p,w,r)$, por lo que es necesariamente convexa en todos los precios.}
\sol{122}{A}{Aplicando el teorema de la envolvente a $\pi(p,w,r) = p q^* - w L^* - r K^*$ se obtiene directamente $\frac{\partial \pi}{\partial p} = q^*$ y $\frac{\partial \pi}{\partial w} = -L^*$.}
\sol{123}{B}{Por convexidad de $\pi$, sus segundas derivadas directas son no negativas: $\frac{\partial q^*}{\partial p} = \frac{\partial^2 \pi}{\partial p^2} \ge 0$ y $\frac{\partial L^*}{\partial w} = -\frac{\partial^2 \pi}{\partial w^2} \le 0$. No existe equivalente a un bien Giffen en demanda de factores.}
\sol{124}{C}{Las condiciones de segundo orden de $\max_{K,L} pf(K,L)-wL-rK$ requieren que el Hessiano de $f$ sea semidefinido negativo (concavidad local), lo que corresponde a rendimientos decrecientes a escala locales y $CMg$ creciente.}
\sol{125}{A}{Maximizar $\pi = pf(K,L)-wL-rK$ respecto a $L$ exige $p \cdot PMg_L = w$ (o $PMg_L = w/p$): la empresa contrata trabajo hasta que el valor monetario del producto adicional iguale su costo salarial.}
\sol{126}{D}{Al no existir restricción de presupuesto fijo, la función de beneficios es convexa en precios y $\frac{\partial L^*}{\partial w} = -\frac{\partial^2 \pi}{\partial w^2} \le 0$. Incluso si un factor es inferior, su demanda incondicional nunca sube al encarecerse.}
\sol{127}{C}{Como $\frac{d CMe}{dq} = \frac{1}{q}(CMg - CMe)$ y $\frac{d CVMe}{dq} = \frac{1}{q}(CMg - CVMe)$, los mínimos ocurren donde $CMg = CMe$ y $CMg = CVMe$ con $CMg$ creciente.}
\sol{128}{B}{Como $CMe = CVMe + CFMe$, en el mínimo de $CVMe$ se tiene $\frac{d CVMe}{dq} = 0$ pero $\frac{d CFMe}{dq} = -CF/q^2 < 0$, por lo que $\frac{d CMe}{dq} < 0$: el $CMe$ sigue disminuyendo y alcanza su mínimo a una escala mayor $q_2 > q_1$.}
\sol{129}{A}{Dado $CV = w L(q)$, derivando respecto a $q$ resulta $CMg = w \frac{dL}{dq} = \frac{w}{PMg_L}$. Análogamente, $CVMe = \frac{w}{PMe_L}$, de modo que el mínimo de $CVMe$ coincide con el máximo de $PMe_L$.}
\sol{130}{D}{Si cierra, su beneficio es $-CF$. Si produce $q>0$, su beneficio es $(p - CVMe)q - CF$. Por tanto, conviene producir siempre que $p \ge \min CVMe$, incluso con pérdidas, porque contribuye a cubrir parte de los costos fijos.}
\sol{131}{B}{Al cerrar se evita tanto $CV(q)$ como $F_c$ (perdiendo solo $CF_h$). Por tanto, operar es mejor que cerrar solo si $pq - CV(q) - F_c \ge 0$, es decir, $p \ge \min \frac{CV(q)+F_c}{q}$.}
\sol{132}{C}{La firma iguala $p = CMg$ (en su tramo creciente por condición de segundo orden) siempre que $p \ge \min CVMe$; para precios menores ofrece $q^* = 0$.}
\sol{133}{D}{Los activos específicos como tuberías subterráneas de una EPS tienen valor de reventa casi nulo fuera de su uso actual (costos hundidos), lo que crea barreras a la salida y riesgos de oportunismo regulatorio (hold-up).}
\sol{134}{A}{Como $CMeLP$ y $CMeCP$ son tangentes, tienen la misma pendiente en el punto de contacto. Si $CMeLP$ es decreciente, $CMeCP$ también debe tener pendiente negativa allí, es decir, opera a la izquierda de su mínimo.}
\sol{135}{C}{Como $C_{LP}(q) = C_{CP}(q; K^*(q))$, por el teorema de la envolvente $\frac{d C_{LP}}{dq} = \frac{\partial C_{CP}}{\partial q}\big|_{K=K^*(q)}$, es decir, $CMgLP$ y $CMgCP$ se cortan exactamente en la producción de tangencia $q_0$.}
\sol{136}{B}{Cuando $\varepsilon_{C,q} = CMgLP/CMeLP < 1$, un aumento del $1\%$ en el producto eleva los costos en menos del $1\%$, por lo que el costo medio de largo plazo desciende (economías de escala).}
\sol{137}{D}{La EME es el menor volumen de producción en el que la firma agota todas las economías de escala disponibles y alcanza el mínimo de su curva de costo medio de largo plazo ($CMeLP$).}
\sol{138}{A}{Al ser $C_{CP}(q) \ge C_{LP}(q)$ con igualdad en $q_0$, la función $C_{CP}(q) - C_{LP}(q)$ tiene un mínimo en $q_0$, por lo que su segunda derivada es no negativa: $\frac{\partial CMgCP}{\partial q} \ge \frac{d CMgLP}{dq}$ (principio de Le Chatelier).}
\sol{139}{B}{Porque en el largo plazo la firma puede ajustar todos sus factores (eliminando la restricción $K=\bar{K}$), su respuesta productiva ante un alza en $p$ es mayor o igual que en el corto plazo: $\frac{\partial q^*_{LP}}{\partial p} \ge \frac{\partial q^*_{CP}}{\partial p}$.}
\sol{140}{C}{Un insumo es inferior si al aumentar la escala de producción $q$ (manteniendo $w, r$ fijos), la firma sustituye hacia tecnologías más intensivas en otros factores y reduce su demanda condicional: $\partial L^c/\partial q < 0$.}
\sol{141}{A}{Por el teorema de Young sobre $C(w,r,q)$, $\frac{\partial CMg}{\partial w} = \frac{\partial^2 C}{\partial w \partial q} = \frac{\partial L^c}{\partial q}$. Si $L$ es inferior ($\partial L^c/\partial q < 0$), un alza en $w$ reduce el $CMg$.}
\sol{142}{D}{Al subir $w$: (1) el efecto sustitución reduce $L^c$ a $q$ fijo ($\le 0$); (2) para un insumo normal, el $CMg$ sube, la firma reduce $q^*$ y el efecto producto reduce aún más $L$. Así, $|\partial L^*/\partial w| \ge |\partial L^c/\partial w|$.}
\sol{143}{B}{Por homogeneidad de grado 0 de $L^c(w,r,q)$, el teorema de Euler implica $w\frac{\partial L^c}{\partial w} + r\frac{\partial L^c}{\partial r} = 0$. Como $\frac{\partial L^c}{\partial w} < 0$, necesariamente $\frac{\partial L^c}{\partial r} > 0$.}
\sol{144}{C}{La subaditividad exige $C(q) < \sum_{i=1}^k C(q_i)$ para toda partición no trivial $\sum q_i = q$. En el caso uniproducto, puede existir subaditividad incluso en un tramo donde $CMeLP$ ya empezó a crecer levemente.}
\sol{145}{D}{Las economías de alcance miden el ahorro de costos derivado de la producción conjunta respecto a la producción especializada separada: $S_C = \frac{C(q_1,0)+C(0,q_2)-C(q_1,q_2)}{C(q_1,q_2)} > 0$.}
\sol{146}{A}{En el caso multiproducto, las economías de escala rayo (o específicas) no bastan para la subaditividad: si existieran deseconomías de alcance, dos firmas especializadas podrían producir más barato que una firma integrada.}
\sol{147}{C}{Las economías de densidad reflejan que, dada una longitud de tuberías instalada, aumentar el volumen servido o el número de conexiones por kilómetro de red diluye los costos fijos de infraestructura y reduce el costo unitario.}
\sol{148}{B}{Si $CMe(q)$ es estrictamente decreciente hasta $q$, entonces $C(q) < \sum C(q_i)$ siempre (suficiente). Sin embargo, una firma puede ser subaditiva en un rango ligeramente posterior al mínimo de $CMe(q)$ (no necesaria).}
\sol{149}{A}{Con $CMg < CMe$, cobrar $p = CMg$ genera ingresos $p q = CMg \cdot q < CMe \cdot q = CT$, provocando pérdidas financieras. Por ello la regulación recurre a tarifas en dos partes, Ramsey-Boiteux o precios iguales a $CMe$.}
\sol{150}{D}{En el progreso neutral de Hicks, $Q = A(t)f(K,L)$: ambos productos marginales crecen en la misma proporción $A(t)$, por lo que $TMST_{L,K} = f_L/f_K$ no varía para una relación $K/L$ dada.}
\sol{151}{C}{El progreso neutral de Harrod es aumentador de trabajo ($F(K, A(t)L)$), mientras que el neutral de Solow es aumentador de capital ($F(A(t)K, L)$) y el de Hicks es $A(t)F(K,L)$.}
\sol{152}{B}{En Cobb-Douglas, $A K^\alpha L^{1-\alpha} = K^\alpha (A^{1/(1-\alpha)}L)^{1-\alpha} = (A^{1/\alpha}K)^\alpha L^{1-\alpha}$, por lo que un avance neutral en cualquiera de los tres sentidos puede reescribirse en los otros dos.}
\sol{153}{D}{Diferenciando logarítmicamente $Q = A f(K,L)$ y usando que en competencia $s_K = \varepsilon_{Q,K}$ y $s_L = \varepsilon_{Q,L}$, se obtiene $\hat{Q} = \hat{A} + s_K \hat{K} + s_L \hat{L}$, de donde $\hat{A} = \hat{Q} - s_K \hat{K} - s_L \hat{L}$.}
\sol{154}{A}{El producto medio es $PMe_L = 18L - L^2$. Derivando e igualando a cero: $18 - 2L = 0 \implies L = 9.0$, donde $PMe_L = PMg_L = 81$.}
\sol{155}{B}{De $PMg_L = 36L - 3L^2 = 0$ resulta $L = 12$. Evaluando en la función de producción: $Q(12) = 18(144) - 1728 = 2592 - 1728 = 864$ unidades.}
\sol{156}{C}{En $L=12$: $PMe_L = 20(12) - 12^2 = 96$ y $PMg_L = 40(12) - 3(12^2) = 48$. Luego $\varepsilon_{Q,L} = \frac{PMg_L}{PMe_L} = \frac{48}{96} = 0.50$.}
\sol{157}{A}{La tasa marginal de sustitución técnica es $TMST_{L,K} = \frac{PMg_L}{PMg_K} = \frac{0.75}{0.25}\frac{K}{L} = 3\left(\frac{16}{8}\right) = 6.0$.}
\sol{158}{D}{Con $\rho = -0.25$, la elasticidad de sustitución técnica es $\sigma = \frac{1}{1 - \rho} = \frac{1}{1 - (-0.25)} = \frac{1}{1.25} = 0.80$.}
\sol{159}{C}{Escalando por $t>0$: $Q(tK, tL) = 2[t^{0.5}(K^{0.5}+L^{0.5})]^3 = t^{1.5}Q(K,L)$. Al ser homogénea de grado $1.5 > 1$, exhibe rendimientos crecientes a escala.}
\sol{160}{B}{En $(2,3)$, $Q = 10$, $PMg_K = L + 2 = 5$ y $PMg_L = K = 2$. Luego $E = \frac{K\cdot PMg_K + L\cdot PMg_L}{Q} = \frac{2(5) + 3(2)}{10} = 1.60$.}
\sol{161}{A}{La tangencia exige $K/L = w/r = 4/16 \implies L = 4K$. En la isocuanta, $\sqrt{4K^2} = 2K = 20 \implies K^* = 10, L^* = 40$. Así, $C = 4(40) + 16(10) = S/\ 320.00$.}
\sol{162}{D}{En el vértice óptimo $2K = 5L = 40$, por lo que $K^* = 20$ y $L^* = 8$. El costo mínimo es $C = 10(8) + 6(20) = S/\ 200.00$.}
\sol{163}{B}{Como $PMg_L/w = 2/6 = 0.333 > PMg_K/r = 3/12 = 0.25$, conviene usar únicamente trabajo: $L^* = 60/2 = 30$, $K^* = 0$, con costo $6(30) = S/\ 180$.}
\sol{164}{C}{La tangencia implica $(K/L)^{0.5} = w/r = 1/4 \implies L = 16K$. Sustituyendo en $Q = 25$: $(K^{0.5} + 4K^{0.5})^2 = 25K = 25 \implies K^* = 1, L^* = 16$, con costo $1(16) + 4(1) = S/\ 20.00$.}
\sol{165}{A}{De $TMST_{L,K} = 2K/L = w/r = 2$ resulta $K = L$. En $Q = K^{0.5}L = L^{1.5} = 8$ se obtiene $L^c = K^c = 4$, y el costo es $8(4) + 4(4) = S/\ 48$.}
\sol{166}{D}{Por el lema de Shephard, $L^c = \frac{\partial C}{\partial w} = q (r/w)^{0.75} = 10(16/1)^{0.75} = 10(8) = 80.0$ unidades.}
\sol{167}{B}{Derivando respecto a $r$: $K^c = \frac{\partial C}{\partial r} = q^2\left(3 + \sqrt{w/r}\right) = 4(3 + 3/2) = 18.0$ unidades.}
\sol{168}{C}{El costo marginal es $CMg = \frac{\partial C}{\partial q} = 6q\sqrt{wr} = 6(5)\sqrt{36} = S/\ 180.00$.}
\sol{169}{A}{Igualando $p\cdot PMg_L = w$: $6\left(\frac{5}{\sqrt{L}}\right) = 15 \implies \sqrt{L} = 2 \implies L^* = 4$ y $Q^* = 20$. El beneficio es $\pi^* = 6(20) - 15(4) - 40 = S/\ 20$.}
\sol{170}{B}{Por simetría ($w=r=2$), $K=L$. La CPO $p\cdot PMg_L = w$ da $8L^{-0.5} = 2 \implies L^* = K^* = 16$, $Q^* = 4(16^{0.5}) = 16$ y $\pi^* = 8(16) - 2(16) - 2(16) = S/\ 64$.}
\sol{171}{D}{Por el lema de Hotelling, $q^*(p,w,r) = \frac{\partial \pi}{\partial p} = \frac{10p}{\sqrt{wr}} = \frac{10(18)}{\sqrt{36}} = 30.0$ unidades.}
\sol{172}{C}{Por el lema de Hotelling, $L^* = -\frac{\partial \pi}{\partial w} = \frac{2.5p^2}{w^{1.5}r^{0.5}} = \frac{2.5(144)}{8 \cdot 3} = 15.0$ unidades.}
\sol{173}{B}{El costo variable medio es $CVMe = 20 - 4q + 0.5q^2$. Minimizando: $-4 + q = 0 \implies q = 4$, donde $p_{\text{cierre}} = \min CVMe = 20 - 16 + 8 = S/\ 12.00$.}
\sol{174}{A}{Igualando $CMg = CMe$: $10 + 2q = \frac{100}{q} + 10 + q \implies q^2 = 100 \implies q = 10$, con $CMe_{\min} = 10 + 10 + 10 = S/\ 30$.}
\sol{175}{C}{El costo medio evitable es $\frac{CV(q)+F_c}{q} = \frac{16}{q} + 4 + q$, cuyo mínimo ocurre en $q = 4$ dando $p_{\min} = 4 + 4 + 4 = S/\ 12.00$.}
\sol{176}{D}{Minimizando respecto a $K$: $4 - q^2/K^2 = 0 \implies K^*(q) = q/2$. Sustituyendo, la envolvente es $C_{LP}(q) = 4q$, que en $q = 30$ vale $S/\ 120.00$.}
\sol{177}{B}{El costo medio es $CMeLP = q^2 - 12q + 60$. Igualando su derivada a cero: $2q - 12 = 0 \implies q = 6$, donde $CMeLP(6) = 36 - 72 + 60 = S/\ 24$.}
\sol{178}{C}{En $q = 10$: $C(10) = 250$, $CMe = 25$ y $CMg = 10 + 10 = 20$. Por tanto, $\varepsilon_{C,q} = \frac{CMg}{CMe} = \frac{20}{25} = 0.80 < 1$, indicando economías de escala.}
\sol{179}{D}{El costo conjunto es $C(30,30) = 400$, mientras que por separado $C(30,0) = 250$ y $C(0,30) = 250$ (suma $500$). Luego $S_C = \frac{500 - 400}{400} = 0.25$.}
\sol{180}{A}{Con un operador, $C(8) = 36 + 64 = S/\ 100$; con dos operadores simétricos, $2C(4) = 2(36 + 16) = S/\ 104$. Al ser $100 < 104$, la función sigue siendo subaditiva aun superando el mínimo de $CMe$.}


\end{document}
